\documentclass[11pt, xcolor=dvipsnames,svgnames,x11names]{article}
\usepackage{xcolor}
\usepackage{xspace}
\usepackage{url}
\usepackage[colorlinks=true,bookmarks=false,linkcolor=black,urlcolor=blue,citecolor=black]{hyperref}
\usepackage{fancyhdr}
\usepackage[top=1in,bottom=1.2in,left=1in,right=1in]{geometry}
\usepackage{multicol}
\usepackage{pgfplots}
\usepackage{tikz}
\usetikzlibrary{circuits.logic.US, patterns}
\usepackage{mathpazo}
\usepackage{biolinum} 
\usepackage[all]{tcolorbox}
\usepackage{../Lectures/rahul_math}
\renewcommand{\familydefault}{\sfdefault}


% --- Custom Color Palette from Image ---
\definecolor{headerRed}{HTML}{A3403D} % Dark red from image top
\definecolor{patternBg}{HTML}{F2D9D7} % Light pinkish bg from image bottom
\definecolor{binaryText}{HTML}{D1A7A4} % Muted color for binary digits

\renewcommand{\texttt}[1]{{\ttfamily\color{headerRed}#1}}


% --- Background Pattern Setup ---
\usepackage[pages=all]{background}
\backgroundsetup{
scale=1,
color=black,
opacity=1,
angle=0,
contents={%
  \begin{tikzpicture}[remember picture,overlay]
    % Top Header Bar
    \fill[headerRed] (current page.north west) rectangle ([yshift=-1.2cm]current page.north east);
    % Bottom Binary Background
    \fill[patternBg] (current page.south west) rectangle ([yshift=3.5cm]current page.south east);
    % Decorative Line (from image)
    \draw[headerRed, line width=2pt] ([xshift=1in, yshift=-2.5cm]current page.north west) -- ([xshift=8in, yshift=-2.5cm]current page.north west);
    % Binary Text Overlay (Sample)
    \node[anchor=south, binaryText, font=\ttfamily\scriptsize, inner sep=0pt, yshift=1.5cm] at (current page.south) {
        \begin{tabular}{c@{\hspace{0.5em}}c@{\hspace{0.5em}}c@{\hspace{0.5em}}c@{\hspace{0.5em}}c@{\hspace{0.5em}}c@{\hspace{0.5em}}c}
        0 0 0 0 1 1 1 1 0 0 0 0 1 1 0 1 1 1 0 0 0 1 \\
        1 0 0 1 0 1 0 0 0 1 0 0 1 1 0 0 1 0 1 1 \\
        0 0 1 1 0 0 1 0 1 1 1 0 1 1 0 1 0 1 0 0 0 \\
        The University of Alabama in Huntsville
        \end{tabular}
    };
  \end{tikzpicture}
}
}

% --- Header Styling ---
\makeatletter
\renewcommand{\maketitle}{%
    \vspace*{0.2cm}
    \begin{center}
        \begin{tcolorbox}[
            enhanced,
            colback=white,
            colframe=headerRed,
            arc=0mm,
            title={\textcolor{white}{\bfseries \@title}},
            coltitle=white,
            fonttitle=\bfseries\large,
            attach boxed title to top left={xshift=10pt, yshift=-\tcboxedtitleheight/2},
            boxed title style={sharp corners, colback=headerRed, frame hidden}
        ]
        \large
        \vspace{10pt}
\noindent
\begin{tabular}{@{}p{3.5cm}p{7cm}r@{}}
\textbf{Student Name:} & \underline{\hspace{7cm}} & \textbf{Points:} 30 \\
\end{tabular}

\vspace{10pt}

\noindent
\begin{tabular}{@{}p{3.5cm}p{7cm}r@{}}
\textbf{A Number:} & \underline{\hspace{7cm}} & \textbf{Duration:} 20 Minutes \\
\end{tabular}
        \end{tcolorbox}
    \end{center}
    \vspace{10pt}
}
\makeatother

\newcommand{\hw}{Quiz 2}
\newcommand{\duedate}{Feb 17, 2026}

\pagestyle{fancy}
\fancyhf{}
\lhead{\small CPE 487/587, Spring 2026}
\rhead{\small \textit{\hw}}
\cfoot{\thepage}
\renewcommand{\headrulewidth}{0pt}

\begin{document}

\title{\hw}
\maketitle

\begin{enumerate}
\item Choose what is true about \texttt{nn.Module}

\begin{enumerate}
\item \texttt{nn.Module} is a function used in PyTorch to train a neural network.
\item \texttt{nn.Module} is an abstract class used to implement a Neural Network Layer.
\item \texttt{nn.Module} is a data type for Neural Networks in PyTorch
\item \texttt{nn.Module} is a neural network parameter.


\end{enumerate}

\item Neural network weights in PyTorch can be of type:
\begin{enumerate}
\item \texttt{nn.Parameter}
\item \texttt{nn.Module}
\item \texttt{nn.Conv2D}
\item \texttt{nn.Weight}
\end{enumerate}

\item Hooks in PyTorch are
\begin{enumerate}
\item functions
\item variables
\item tensors
\item classes
\end{enumerate}


\item What is the difference between parameters and buffers in the realm of PyTorch  \texttt{nn.Module}?

\begin{tcolorbox}[colback=white, colframe=headerRed!75!black]
\vspace{1.5em}
\noindent\rule{\textwidth}{0.4pt}\\[1.5em]
\noindent\rule{\textwidth}{0.4pt}\\[1.5em]
\noindent\rule{\textwidth}{0.4pt}\\[1.5em]
\noindent\rule{\textwidth}{0.4pt}
\vspace{0.5em}
\end{tcolorbox}


\item What happens if you don't register a member variable tensor as parameters in PyTorch \texttt{nn.Module}?

\begin{tcolorbox}[colback=white, colframe=headerRed!75!black]
\vspace{1.5em}
\noindent\rule{\textwidth}{0.4pt}\\[1.5em]
\noindent\rule{\textwidth}{0.4pt}\\[1.5em]
\noindent\rule{\textwidth}{0.4pt}\\[1.5em]
\noindent\rule{\textwidth}{0.4pt}
\vspace{0.5em}
\end{tcolorbox}


\item When are forward hooks executed?

\begin{enumerate}
\item The hook are called every time after \texttt{backward()}.
\item The hook are called just before calling \texttt{forward()}.
\item The hook are called every time after weights are updated..
\item The hook are called every time after \texttt{forward()}.
\end{enumerate}

\item Is it possible to implement a neural network in PyToch without inhering from \texttt{nn.Module}?

\begin{enumerate}
\item Yes
\item No
\end{enumerate}


\item Consider a Network definition 

          \begin{minted}[
                linenos,
                frame=lines,
                framesep=2mm,
                bgcolor=pink!10,
                fontsize=\footnotesize,
                breaklines
            ]{Python}
class Network(nn.Module):
	def __init__(self):
		super().__init__()
		self.layer1 = nn.Linear(10, 20)
		self.layer2 = nn.Linear(20, 5)
	def forward(self, x):
		return self.layer2(self.layer1(x))
\end{minted}

How many trainable parameters are there in the above network (including biases)?

\begin{enumerate}
\item 325
\item 300
\item 55
\item 400
\end{enumerate}

\item Imagine a dataset consisting of PNG images with four channels: RGBA. You decided to use a fully connected neural network for creating an image classfier. There 10000 images in your training dataset. Each image has dimension $200\times 200$. How many features will be there in the input feature matrix to the neural network?

\begin{enumerate}
\item 10000
\item 1600000000
\item 40000
\item 400000000
\end{enumerate}

\item Minpooling downsamples the image data.
\begin{enumerate}
\item True
\item False
\end{enumerate}

\item An example of order-5 Tensor:
\begin{tcolorbox}[colback=white, colframe=headerRed!75!black]
\vspace{1.5em}
\noindent\rule{\textwidth}{0.4pt}\\[1.5em]
\noindent\rule{\textwidth}{0.4pt}
\vspace{0.5em}
\end{tcolorbox}

\item Consider a convolution layer definition


          \begin{minted}[
                linenos,
                frame=lines,
                framesep=2mm,
                bgcolor=pink!10,
                fontsize=\footnotesize,
                breaklines
            ]{Python}
conv = nn.Conv2d(in_channels=4,  
                    out_channels=6, 
                    kernel_size=3, 
                    stride=2,
                    padding=1)
\end{minted}

What is the dimension of an example tensor that can serve as an input to the above convolution layer?

\begin{enumerate}
\item \texttt{torch.Size([2, 1, 5, 6])]}
\item \texttt{torch.Size([2, 6, 32, 32])]}
\item \texttt{torch.Size([2, 3, 5, 6])]}
\item \texttt{torch.Size([2, 4, 32, 32])]}
\end{enumerate}

\item For the convolution operation in Q12, how many filters will be learnt upon training?
\begin{enumerate}
\item 4
\item 3
\item 6
\item 2
\end{enumerate}

\item If you are performing regression with a neural network, what should be the final layer in a neural network?

\begin{enumerate}
\item A linear layer
\item A sigmoid layer
\item A relu layer
\item A tanh layer
\end{enumerate}

\item Applying max pooling with a kernel size of $3\times 3$ to the image of size $5\times 5$ with a stride of $2$, the resulting output has size 

\begin{enumerate}
\item $2\times 2$
\item $3\times 3$
\item $5\times 5$
\item $3\times 5$
\end{enumerate}

\end{enumerate}

\end{document}